\section{Model development}
The smoothed transition table uses the last observed event. The current-row logistic regression uses one-hot category encoding and standardized numeric values, with imputation fitted on the fitting split. These baselines help measure the gain from longer histories and nonlinear features.

\subsection{Classical model}
We tune LightGBM \cite{lightgbm} on the feature groups in Table~1 with the categorical event pair excluded. The earlier ablation on these validation machines motivated this choice. The final test set remains held out. At the selected tree settings, we restore the 196-category pair in one ablation and remove the individual lags in another. We also fit on four sampled cutoffs per machine to check whether using fewer overlapping training examples changes performance. Categorical vocabularies are fitted on fitting prefixes and unknown levels are encoded as missing for prediction. Trees handle missing numerical values directly. The model is trained from scratch using the multiclass objective.

\subsection{First deep model}
% First deep model (GRU). Included from 02_models.tex.
The recurrent model reads the prefix event by event instead of a fixed summary. Each step is the concatenation of a 32-dimensional embedding of the event class, 8-dimensional embeddings of subsystem, severity, shift, operator group, device family and firmware channel, and numeric inputs: the log time gap to the previous event, the position in the log, all 16 measurements (including the four telemetry channels the classical sensor history omits) and a missing-value flag for each. Measurements are standardized with means and deviations from the fitting machines, and missing values are set to zero so that the flag carries the information. A linear layer with ReLU maps each step to $h$ units, a unidirectional GRU or LSTM processes the sequence, and layer normalization and a linear head give 14 logits at every step.

The recurrence only moves forward in time, so the output at step $t$ depends on events $0,\dots,t$ and is the prediction for a prefix of length $t+1$. We use this in two ways. First, training is many-to-many: one pass over a fitting machine supervises every step whose prefix length lies in 90--250, giving about 95 targets per machine per epoch instead of 12 sampled cutoffs. The classical cutoff ablation, where four instead of twelve cutoffs per machine raised log loss by 0.035, suggests that more training prefixes help. Second, variable prefix lengths need no packing. Shorter sequences are padded on the right, padding cannot affect earlier outputs, and we read the output at the last observed event. For the short-context ablation the model instead sees a random chunk of $K+32$ events and is trained only on steps with at least $K$ events of history, matching the $K$ events it receives at evaluation.

Training minimizes cross-entropy, which equals log loss, with AdamW (weight decay $10^{-4}$), batches of 64 machines and gradient-norm clipping at 1. The learning rate halves when the early-stop loss stops improving, and training stops after four epochs without improvement or after 40 epochs. We keep the weights from the best epoch. The early-stop loss is computed on the 12 manifest cutoffs of the 520 early-stop machines, so it mirrors the validation task. We then fit one temperature $T$ on the 520 calibration machines by minimizing the log loss of $\mathrm{softmax}(z/T)$. Dividing all logits by $T$ preserves their order, so accuracy is unchanged and only confidence is adjusted.


\subsection{Second deep model}
\todo{Person C. Document the architecture, temporal context, training settings and comparison to the first deep model.}
